\documentclass[book.tex]{subfiles}
\begin{document}

\chapter{SELFIES Representations}

\label{chap:selfies}
\noindent\rule{11cm}{0.4pt} \\
\textbf{Prerequisites:} \autoref{chap:graphconvs}, \autoref{chap:equivariant_networks}, \autoref{chap:molecular_dynamics},  \autoref{chap:dft} \\
\textbf{Difficulty Level:} ** \\
\noindent\rule{11cm}{0.4pt}
\vspace{5mm}

"I will talk about SELFIES, a 100\% robust molecular string representation. The main feature is that every combination of letters from the alphabet gives a meaningful molecule. The main application for SELFIES are (deep) generative models, for the design of new functional molecules. In contrast to other representations, here every output of the model is syntactically and semantically correct thus no posterior filtering or adaptations of the models are necessary. I will show how this allowed, for example, the direct application of curiosity-based exploration as well as DeepDreaming techniques. Finally, I will highlight two works by other groups that have shown that SELFIES outperform all other representations in string2string and img2string translation tasks. Thus in some sense, it indicates that SELFIES is easier to "understand" for machines -- which raises several exciting questions."

References \cite{thiede2020curiosity}, \cite{shen2021deep}, \cite{krenn2020self}

\end{document}